\documentclass[11pt,a4paper]{article}

\usepackage[utf8]{inputenc}
\usepackage[T1]{fontenc}
\usepackage[margin=1in]{geometry}
\usepackage{booktabs}
\usepackage{hyperref}
\usepackage{amsmath}
\usepackage{graphicx}
\usepackage{xcolor}
\usepackage{microtype}
\usepackage{parskip}
\usepackage{natbib}
\usepackage{listings}
\usepackage{caption}
\usepackage{multirow}
\usepackage{float}

\hypersetup{
  colorlinks=true,
  linkcolor=blue!60!black,
  citecolor=blue!60!black,
  urlcolor=blue!60!black,
}

\lstset{
  basicstyle=\ttfamily\small,
  breaklines=true,
  frame=single,
  backgroundcolor=\color{gray!8},
}

\title{\textbf{Near-Perfect Session Retrieval on LongMemEval:}\\
       \large MemoryKG's Hybrid Semantic-Graph Architecture\\
       Achieves 99.4\% Recall@10 and 100\% Recall@30 Across 500 Questions}

\author{Eric G. Suchanek, PhD\\
        \texttt{suchanek@flux-frontiers.com}}

\date{April 25, 2026}

% ============================================================
\begin{document}
\maketitle

% ============================================================
\begin{abstract}
We evaluate MemoryKG, a conversational memory system based on
DocKG~\citep{dockg2026}, on the LongMemEval
benchmark~\citep{wang2024longmemeval}, a dataset of 500 questions spanning six
question types that require locating specific sessions within a multi-session
conversation corpus of up to 500 sessions.
MemoryKG achieves \textbf{99.4\% session-level Recall@10} and
\textbf{100\% Recall@30} across all question types, with an NDCG@10 of 0.943,
on the full 500-question test set.
Perfect recall is attained at $k{=}30$ with no LLM required at any stage ---
neither for indexing nor for retrieval.
The hardest question types are \textit{temporal-reasoning} (Recall@10 = 98.5\%)
and \textit{multi-session} (Recall@10 = 99.2\%), where evidence spans multiple
chronologically ordered sessions.
The easiest types --- \textit{single-session-user}, \textit{single-session-assistant},
and \textit{single-session-preference} --- each reach 100\% Recall@10.
Graph expansion through structured \textsc{Contains}, \textsc{Next}, and
\textsc{References} edges closes the gap between semantic seeding and
complete evidence recovery, requiring only $k{=}30$ seeds to exhaust all 500
questions.
A single shared sentence-transformer embedder (BAAI/bge-small-en-v1.5)
processes each per-question corpus in approximately 0.36 seconds on Apple Silicon.
\end{abstract}

% ============================================================
\section{Introduction}
\label{sec:intro}

Long-term conversational memory retrieval --- locating the right session from
a large multi-session corpus given a natural-language question --- is
increasingly important for AI assistants operating over months of interaction history.
Prior approaches include full-context LLM inference, structured memory extraction
\citep{mem0}, and flat dense retrieval over session text.
These approaches differ in accuracy, latency, cost, and privacy guarantees.

The LongMemEval benchmark \citep{wang2024longmemeval} provides a structured
evaluation across six question types, each requiring the system to locate
the correct session(s) from a corpus of up to 500 long-horizon conversation
sessions.
The benchmark distinguishes single-session retrieval (facts, preferences,
assistant claims) from multi-session and temporally-ordered evidence.

MemoryKG~\citep{memorykg2026} builds a lightweight knowledge graph over verbatim
conversation text, where semantic vector search seeds initial candidates and
structured graph traversal expands the result set.
MemoryKG is a direct specialisation of DocKG~\citep{dockg2026}, a hybrid
semantic-graph system for general document corpora.
In this paper we report MemoryKG's performance on LongMemEval and analyse the
retrieval characteristics of each question type.

% ============================================================
\section{The LongMemEval Benchmark}
\label{sec:benchmark}

LongMemEval \citep{wang2024longmemeval} is a benchmark for evaluating chat assistants
on long-term interactive memory.
It contains 500 questions drawn from realistic multi-session conversation histories.
Each question is paired with one or more \emph{evidence sessions} --- the specific
conversation turns or sessions whose content is necessary to answer the question
correctly.

The six question types are:

\begin{description}
  \item[Single-Session User ($n{=}70$)] Questions about facts the user stated in a
        single session.
  \item[Single-Session Assistant ($n{=}56$)] Questions about claims or information
        the assistant provided in a single session.
  \item[Single-Session Preference ($n{=}30$)] Questions about user preferences expressed
        within a single session.
  \item[Knowledge Update ($n{=}78$)] Facts that the user updated or corrected across
        sessions; the system must retrieve the most recent statement.
  \item[Multi-Session ($n{=}133$)] Questions whose answer requires aggregating
        information from two or more separate sessions.
  \item[Temporal Reasoning ($n{=}133$)] Questions that require reasoning about the
        order or timing of events across sessions.
\end{description}

We use \emph{session-level retrieval recall} as our primary metric: the fraction
of evidence sessions whose text is found in the top-$k$ retrieved nodes.
We also report NDCG@$k$ (Normalised Discounted Cumulative Gain), which rewards
placing the correct session higher in the ranked list.
These are binary retrieval metrics, not end-to-end QA accuracy.
The LongMemEval paper reports end-to-end QA accuracy (requiring an LLM to generate
an answer); we report whether the evidence session is present in the retrieved context.
These are complementary but distinct measures.

% ============================================================
\section{MemoryKG Architecture}
\label{sec:arch}

MemoryKG represents a conversation corpus as a four-layer knowledge graph:

\begin{enumerate}
  \item \textbf{Document layer.} Each conversation session is a root document node
        containing the full session text.
  \item \textbf{Section layer.} Markdown headings (one per speaker turn) produce
        section nodes connected to their parent document by \textsc{Contains} edges.
  \item \textbf{Chunk layer.} Section content is split into sentence-aligned chunks;
        adjacent chunks are linked by \textsc{Next} edges.
  \item \textbf{Semantic index.} All nodes are embedded with a sentence transformer
        and stored in LanceDB for approximate nearest-neighbour retrieval.
\end{enumerate}

\subsection{Query pipeline}

A query $q$ is processed in two stages:

\paragraph{Stage 1 --- Semantic seeding.}
The top-$k$ nodes by cosine similarity to the query embedding are retrieved from LanceDB.
These form the \emph{seed set} $S$.

\paragraph{Stage 2 --- Graph expansion.}
Starting from $S$, BFS expansion walks outward through the graph for $h$ hops along
any edge in the relation set $\mathcal{R}$.
The default relation set is
\begin{equation*}
  \mathcal{R} = \{\textsc{Contains},\;\textsc{Next},\;\textsc{References},\;
    \textsc{Similar\_To},\;\textsc{Has\_Topic},\;\textsc{Mentions\_Entity},\;
    \textsc{Has\_Keyword}\}.
\end{equation*}
Nodes reachable within $h$ hops of any seed are added to the candidate pool and ranked
by a composite key combining seed distance, hop depth, and structural position.

For LongMemEval evaluations we use $k{=}10$ seeds and $h{=}1$ hop.

\subsection{Why graph expansion helps for long-horizon retrieval}

Consider a Knowledge Update question: ``What is the user's current job title?''
The earliest session might read ``I just started as a junior analyst,'' while a
later session reads ``Got promoted to senior analyst last month.''
A flat vector search that ranks the earlier session higher will return a stale
answer.
Graph expansion helps via \textsc{Contains} edges: the parent document node for
each session contains the full session text, and temporal ordering is encoded in
\textsc{Next} edges between sessions.
When the correct (most recent) session is seeded by any related chunk, expansion
propagates recall upward to the session level.

For temporal reasoning questions, the same mechanism resolves ordering ambiguity:
if any chunk from the target session appears in the seed set, traversal lifts the
entire session into the candidate pool regardless of whether the question vocabulary
matches the session's surface text.

% ============================================================
\section{Experimental Setup}
\label{sec:setup}

\subsection{Corpus format}

For each question, all sessions in the user's history are written to a temporary
directory as Markdown files (one file per session).
Each speaker turn becomes a \texttt{\#\#}-headed section:

\begin{lstlisting}
# Session 2024-03-15

## User

I just got promoted to senior analyst last month.

## Assistant

Congratulations! That's a big step --- what changed in your role?

## User

More client-facing work, mostly. I'm leading the quarterly reviews now.
\end{lstlisting}

The heading chunk strategy splits at \texttt{\#\#} boundaries, so each turn
becomes one chunk node.

\subsection{Build configuration}

\begin{itemize}
  \item Chunk strategy: \texttt{heading} (no per-file embedding at parse time)
  \item Topic / entity / keyword enrichment: disabled (speed)
  \item SIMILAR\_TO edge discovery: disabled (speed)
  \item Workers: 1 (per-item builds are already fast)
  \item Batch size: 512 embeddings
\end{itemize}

\subsection{Shared embedder}

A single \texttt{SentenceTransformerEmbedder}
(\texttt{BAAI/bge-small-en-v1.5}, 384-dimensional)
is initialised once and reused across all questions.
This avoids the per-question model reload overhead and is the primary contributor
to the 0.36 s/question throughput.
The full 500-question evaluation completes in 181.6 seconds on an Apple M5 Max.

\subsection{Recall and NDCG computation}

For each question, let $E = \{e_1, \ldots, e_n\}$ be the set of evidence session texts and
$N = \{n_1, \ldots, n_m\}$ be the ranked list of retrieved node texts.
Session-level recall at $k$ is:
\[
  \text{Recall@}k = \frac{1}{|E|}
    \sum_{e \in E} \mathbf{1}\!\left[\exists\, n \in N_k :
      e \subset n \;\vee\; n \subset e
    \right]
\]
where $\subset$ denotes case-insensitive substring containment and $N_k$ is the
top-$k$ subset of $N$.
NDCG@$k$ is computed over binary relevance labels using the standard formulation:
\[
  \text{NDCG@}k = \frac{\text{DCG@}k}{\text{IDCG@}k}, \qquad
  \text{DCG@}k = \sum_{i=1}^{k} \frac{\text{rel}_i}{\log_2(i+1)}
\]
where $\text{rel}_i \in \{0,1\}$ indicates whether position $i$ in the ranking
contains an evidence session.

% ============================================================
\section{Results}
\label{sec:results}

\subsection{Session-level retrieval across $k$}

Table~\ref{tab:recall_at_k} shows MemoryKG's session-level Recall@$k$ and NDCG@$k$
across increasing retrieval depths.

\begin{table}[H]
\centering
\caption{MemoryKG session-level retrieval metrics on LongMemEval (500 questions).}
\label{tab:recall_at_k}
\begin{tabular}{rrr}
  \toprule
  \textbf{$k$} & \textbf{Recall@$k$} & \textbf{NDCG@$k$} \\
  \midrule
   1  & 0.904 & 0.904 \\
   3  & 0.974 & 0.942 \\
   5  & 0.984 & 0.942 \\
  10  & 0.994 & 0.943 \\
  30  & \textbf{1.000} & 0.942 \\
  50  & 1.000 & 0.942 \\
  \bottomrule
\end{tabular}
\end{table}

Perfect recall is reached at $k{=}30$: all 500 evidence sessions appear in the
top-30 retrieved nodes.
The NDCG@$k$ stabilises at 0.942--0.943 for $k \geq 3$, indicating that the
correct sessions consistently rank near the top of the list once they appear.

\subsection{Per-type breakdown}

Table~\ref{tab:type_results} shows Recall@5, Recall@10, and NDCG@10 for each
LongMemEval question type.

\begin{table}[H]
\centering
\caption{MemoryKG retrieval metrics on LongMemEval by question type (top-10, hop=1).}
\label{tab:type_results}
\begin{tabular}{lrrrr}
  \toprule
  \textbf{Question Type} & \textbf{$n$} & \textbf{Recall@5} & \textbf{Recall@10} & \textbf{NDCG@10} \\
  \midrule
  knowledge-update          &  78 & 1.000 & 1.000 & 0.984 \\
  multi-session             & 133 & 0.992 & 0.992 & 0.949 \\
  single-session-assistant  &  56 & 1.000 & 1.000 & 1.000 \\
  single-session-preference &  30 & 1.000 & 1.000 & 0.913 \\
  single-session-user       &  70 & 0.986 & 1.000 & 0.932 \\
  temporal-reasoning        & 133 & 0.955 & 0.985 & 0.901 \\
  \midrule
  \textbf{Overall}          & \textbf{500} & \textbf{0.984} & \textbf{0.994} & \textbf{0.943} \\
  \bottomrule
\end{tabular}
\end{table}

Four of the six question types achieve 100\% Recall@10.
The two types that do not --- \textit{temporal-reasoning} (98.5\%) and
\textit{multi-session} (99.2\%) --- are those requiring evidence from multiple
or temporally ordered sessions.

\subsection{Comparison with literature baselines}

Table~\ref{tab:lit_comparison} places MemoryKG in the context of systems
evaluated in the LongMemEval literature~\citep{wang2024longmemeval}.
Those results are reported as end-to-end QA accuracy, which requires an LLM to
generate a free-form answer given retrieved context.
MemoryKG's retrieval recall is a strictly easier task: it measures whether the
evidence is \emph{present} in the retrieved context, not whether a reader can
extract the correct answer from it.
The two metrics are complementary but not directly comparable; the table is
included for broader orientation.

\begin{table}[H]
\centering
\caption{Context comparison.
         \dag~QA accuracy from~\citet{wang2024longmemeval} (requires LLM answer
         generation); ours is retrieval recall (no LLM). Different metrics;
         not directly comparable.}
\label{tab:lit_comparison}
\begin{tabular}{lrll}
  \toprule
  \textbf{System} & \textbf{Score} & \textbf{Metric} & \textbf{LLM required} \\
  \midrule
  MemoryKG (ours, R@10)       & \textbf{99.4\%} & Retrieval recall & None \\
  MemoryKG (ours, R@30)       & \textbf{100\%}  & Retrieval recall & None \\
  \midrule
  GPT-4o full-context\dag     & $\sim$66\%      & QA accuracy      & Yes  \\
  GPT-4o + RAG\dag            & $\sim$50\%      & QA accuracy      & Yes  \\
  GPT-4-Turbo full-context\dag& $\sim$56\%      & QA accuracy      & Yes  \\
  Claude-3.5-Sonnet\dag       & $\sim$52\%      & QA accuracy      & Yes  \\
  ReadAgent\dag               & $\sim$42\%      & QA accuracy      & Yes  \\
  Zero-shot (no memory)\dag   & $\sim$25\%      & QA accuracy      & Yes  \\
  \bottomrule
\end{tabular}
\end{table}

\subsection{Miss analysis}

At $k{=}10$, exactly 3 out of 500 questions returned no evidence session in the
top-10 retrieved nodes.
These three misses correspond to question IDs \texttt{10d9b85a},
\texttt{gpt4\_468eb064}, and \texttt{eac54add}.
All three were recovered by $k{=}30$, confirming that the evidence exists in the
index but ranks below position 10 in these cases.
The miss rate at $k{=}10$ is 0.6\% (3/500).

% ============================================================
\section{Analysis}
\label{sec:analysis}

\subsection{Why temporal-reasoning is the hardest type}

Temporal-reasoning questions ask about the ordering or recency of events across
sessions (e.g., ``When did the user first mention they were learning Spanish?'').
The evidence session may contain only a passing reference to the target event,
while semantically closer sessions discuss unrelated topics.
Vocabulary mismatch between the question and the target session is highest in
this category: a question about ``first mention'' may correspond to a session
that simply reads ``I started a Spanish class'' without temporal markers.
Graph expansion partially compensates --- if any chunk from the same session
appears in the seed set, \textsc{Next} and \textsc{Contains} edges recover the
full session --- but 2 out of 133 temporal-reasoning questions (1.5\%) still
miss at $k{=}10$.

\subsection{Why multi-session questions remain hard}

Multi-session questions require evidence from two or more sessions simultaneously.
At $k{=}10$, the system must place all required sessions within the top-10 result,
which becomes more constrained as the number of required sessions grows.
Unlike the ConvoMem tier model~\citep{pakhomov2025convomem}
(which fixes the number of evidence messages),
LongMemEval multi-session questions vary in evidence count.
The 1 miss at $k{=}10$ for this category represents a case where one required
session ranked below position 10 while the other(s) ranked highly.

\subsection{The NDCG plateau and its implications}

NDCG@$k$ rises sharply from 0.904 at $k{=}1$ to 0.942 at $k{=}3$ and then
plateaus at 0.942--0.943 for all larger $k$.
This plateau indicates that additional retrieved nodes beyond $k{=}3$ do not
substantially improve the rank of correct sessions --- the correct sessions are
already near the top of the list when they are present.
The gap between Recall@$k$ and the NDCG plateau reflects the 3 hard misses
that only appear beyond position 10.
For a downstream reader consuming the top-$k$ retrieved context, $k{=}10$ is
effectively optimal: additional items improve coverage without meaningfully
improving precision.

\subsection{Throughput and scaling}

At 0.36 s/question on Apple Silicon (M5 Max), MemoryKG processes 500 questions
in 181.6 seconds.
The dominant cost is LanceDB indexing (Phase 2) per question, which scales with
the number of sessions in the corpus.
LongMemEval corpora are larger than ConvoMem corpora (up to 500 sessions vs.\ up
to 300 conversations), accounting for the higher per-question time (0.36 s vs.\
0.07 s on ConvoMem).
The shared embedder eliminates model reload overhead; the remaining cost is
proportional to corpus size and is bounded by MPS-accelerated LanceDB ingestion.

% ============================================================
\section{Conclusion}
\label{sec:conclusion}

MemoryKG achieves 99.4\% Recall@10 and 100\% Recall@30 on the 500-question
LongMemEval benchmark, with an NDCG@10 of 0.943 and no LLM required at any stage.
Four of six question types reach perfect Recall@10.
The two remaining types --- temporal-reasoning and multi-session --- reflect the
fundamental difficulty of retrieving evidence from multi-session contexts with
temporal ordering requirements; both approach but do not reach 100\% at $k{=}10$,
and both are fully resolved by $k{=}30$.

The hybrid architecture --- semantic seeding combined with structured graph
expansion --- enables high-precision retrieval without the latency or cost of
LLM inference.
The NDCG plateau at 0.942 from $k{=}3$ onward confirms that when evidence
sessions are retrieved, they rank near the top of the result list, making the
system suitable for downstream readers with tight context budgets.

\paragraph{Limitations.}
Recall measures whether the evidence session is \emph{present} in the retrieved
context; it does not measure whether a downstream reader can correctly extract
the answer.
End-to-end QA accuracy on LongMemEval with a MemoryKG retrieval front-end is
left for future work.
The 3 misses at $k{=}10$ are all resolved by $k{=}30$ and are concentrated in
the temporal-reasoning category, suggesting that temporal ordering cues are the
primary remaining failure mode.

\paragraph{Reproducibility.}
All benchmark scripts, result JSON files, and this writeup are included in the
MemoryKG repository.
The full evaluation can be reproduced with:
\begin{lstlisting}
python benchmarks/longmemeval/longmemeval_bench.py --limit 500
\end{lstlisting}

% ============================================================
\bibliographystyle{plainnat}
\begin{thebibliography}{9}

\bibitem[Wang et al.(2024)]{wang2024longmemeval}
Wang, X., et al. (2024).
\newblock {LongMemEval: Benchmarking Chat Assistants on Long-Term Interactive
  Memory}.
\newblock \textit{arXiv preprint arXiv:2410.10813}.
\newblock \url{https://arxiv.org/abs/2410.10813}

\bibitem[Suchanek(2026a)]{memorykg2026}
Suchanek, E.~G. (2026).
\newblock {MemoryKG: A Hybrid Semantic-Graph Knowledge Base for Conversational Memory}.
\newblock \url{https://github.com/Flux-Frontiers/memory_kg}

\bibitem[Suchanek(2026b)]{dockg2026}
Suchanek, E.~G. (2026).
\newblock {DocKG: Semantic Knowledge Graph for Document Corpora}.
\newblock Version 0.11.0. Flux-Frontiers. Initial commit: 2026-03-08.
\newblock DOI: \url{https://doi.org/10.5281/zenodo.19770973}

\bibitem[Mem0(2024)]{mem0}
{Mem0 Team}. (2024).
\newblock {Mem0: The Memory Layer for Personalized AI}.
\newblock \url{https://mem0.ai}

\bibitem[Pakhomov et al.(2025)]{pakhomov2025convomem}
Pakhomov, E., Nijkamp, E., \& Xiong, C. (2025).
\newblock {ConvoMem Benchmark: Why Your First 150 Conversations Don't Need RAG}.
\newblock \textit{arXiv preprint arXiv:2511.10523}.
\newblock \url{https://arxiv.org/abs/2511.10523}

\end{thebibliography}

\end{document}
